
\section{Conclusion}
In this study, we hypothesized that Timeboost exhibits a lower MEV capture rate under volatile market conditions and empirically tested this hypothesis. Our methodology uses realized values to construct a noisy proxy for an unobserved variable that captures bidders' forecasts of future volatility. We developed a bidder valuation model and proposed using the coefficient on the variance of integrated variance ($\Var(\IV)$) as the risk premium for the uncertainty in the value of access to the express lane (EL). 
Subsequently, we compared the regression results of our model with those of a baseline model that excluded a risk premium, demonstrating improvements in fit and a reduction in downward bias in the restricted model. We further tested our hypothesis under various conditions to check the robustness of the results.

This study shows that variance risk premiums are statistically significant and negative. Specifically, it demonstrates that the full model provides a better explanation of the findings compared to the restricted model. As expected, the restricted model yields a downward-biased slope relative to the full model, reinforcing the overall argument. We further tested our hypothesis under various conditions, checking the results with subsamples, using alternative measures of volatility, dropping the rational-expectations assumption, and using forecasted rather than realized values to assess robustness.

Based on the results, we provided an outline to improve the current Timeboost auction implementation and increase the sequencer's revenue. Specifically, we offered two suggestions: a more advanced method for setting the reserve price and extending the round length, which helps bidders better estimate future (averaged) volatility, leading to less discounting. We also briefly mentioned the option of switching to a subscription model, which is already popular in traditional finance (TradFi).

Our methodology exhibits several notable limitations. Primarily, the model's simplicity results in the omission of numerous real-world complexities. For example, the model assumes that the liquidity available around the CEX price remains constant throughout the analysis period and restricts its scope to ETH as the only volatile asset within AMMs. Furthermore, reliance on the CIPV framework may not fully capture the intricacies of actual market dynamics. The use of future realized variance and Newey-West estimators, which are inherently unknown at the time of bidding, as noisy proxies for bidders’ forecasts, introduces attenuation bias. Finally, the substantial rate of data censorship poses significant challenges to the robustness and validity of the inferences drawn, despite the use of tobit regression.

Despite of aforementioned limitations, this study is one of the first empirical investigations of Timeboost, examining the determinants of bidders' valuations and the efficacy of the express lane auctions. As such, it provides a foundational reference for subsequent research that employs more advanced methodologies and incorporates realistic assumptions. Ultimately, this study aims to enhance the understanding of time-based sequencing mechanisms within auction frameworks.

% \begin{table}[htbp]
% \centering
% \caption{Heteroskedastic Tobit Regression Results (with 1/P): Bidder 0x8c6f}
% \label{tab:tobit_inv_p_8c6f}
% \begin{tabular}{lcccc}
% \toprule
%  & \multicolumn{2}{c}{Linear (t-dist)} & \multicolumn{2}{c}{Log-Log (Gaussian)} \\
% \cmidrule(lr){2-3} \cmidrule(lr){4-5}
%  & Reduced & Full & Reduced & Full \\
% \midrule
% \multicolumn{5}{l}{\textit{Panel A: Location Coefficients}} \\[3pt]
% Intercept & -0.6861 (0.0142) & -0.7161 (0.0142) & 52.8063 (0.0584) & 51.2007 (0.0579) \\
% $1/P$ & 0.7530 (0.0049) & 0.7392 (0.0049) & 1.0996 (0.0069) & 0.8142 (0.0071) \\
% $\mathbb{E}[\text{IV}]/\sqrt{P}$ & 0.1822 (0.0012) & 0.2109 (0.0010) & 0.4220 (0.0013) & 1.0195 (0.0050) \\
% $\text{Var}(\text{IV})$ & -- & -0.9150 (0.0057) & -- & -0.3348 (0.0027) \\
% \midrule
% \multicolumn{5}{l}{\textit{Panel B: Scale Coefficients (log-scale)}} \\[3pt]
% Scale Intercept & 0.3188 & 0.3402 & -0.0582 & -0.0879 \\
% \midrule
% \multicolumn{5}{l}{\textit{Panel C: Model Fit}} \\[3pt]
% Log-Likelihood & -748067.73 & -745897.86 & -346747.51 & -339103.33 \\
% AIC & 1496145.45 & 1491807.71 & 693503.02 & 678216.67 \\
% BIC & 1496197.89 & 1491870.63 & 693544.97 & 678269.10 \\
% Degrees of Freedom (t) & 0.78 & 0.80 & $\infty$ & $\infty$ \\
% N (total) & 264959 & 264959 & 264959 & 264959 \\
% N (uncensored) & 233320 & 233320 & 233320 & 233320 \\
% \midrule
% \multicolumn{5}{l}{\textit{Panel D: LR Test (Reduced vs Full)}} \\[3pt]
% LR Statistic & \multicolumn{2}{c}{4339.74} & \multicolumn{2}{c}{15288.35} \\
% $p$-value & \multicolumn{2}{c}{$<$0.001***} & \multicolumn{2}{c}{$<$0.001***} \\
% \bottomrule
% \end{tabular}
% \begin{tablenotes}[flushleft]
% \footnotesize
% \item \textit{Notes:} Standard errors in parentheses. Linear models use Student's t-distributed residuals; log-log models use Gaussian residuals. Left-censoring at 0.001 ETH. *** $p<0.001$, ** $p<0.01$, * $p<0.05$.
% \end{tablenotes}
% \end{table}
